% ---
% id: meat-core/frank-wolfe/lem:certified-bracket
% kind: lem
% title: "The certified bracket from enclosures of the gradient"
% summary: >-
%   At a point with positive definite blocks and positive probabilities the
%   gradient paired with the point is the objective, so the bracket is the sum
%   of its two lower bounds and needs only enclosures of the gradient.
% prominence: standard
% status: implemented
% uses: [thm:firstorder, lem:feasible-side-dual, lem:fo-khyp, def:first-order-setting, def:relaxed-feasible-set, cor:fo-dropped, meat-core/lem:objective-homogeneous, meat-core/cor:full-gradient, meat-core/thm:gradient, meat-core/lem:branch-support, meat-core/thm:branch-formula, meat-core/thm:choi, meat-core/def:declaration, meat-core/def:admissible-declaration, meat-core/def:admissible-acceptance-set, meat-core/def:joint-sectors, meat-core/lem:joint-sectors, meat-core/def:keyinert, meat-core/def:keyactive, meat-core/def:dropped, meat-core/prop:drop-bound, foundations/def:kl, foundations/cor:least-eigenvalue-over-enclosure]
% mentions: [lem:certified-bracket-enclosed, cor:certified-keylength, req:certified-evaluation]
% assumes: []
% refs: []
% provenance: newly-derived
% implements: certiqs_key.solvers.frank_wolfe.certified:certified_bracket
% lean:
% ---

Part~(iv) of Theorem~\ref{thm:firstorder} lets lower bounds on the two linear
minima stand in for the pairings with an oracle's return. The objective is
positively homogeneous of degree one
(Lemma~\ref{lem:objective-homogeneous}), so its derivative along the point
itself is its value there, and the bracket loses the objective and the two
pairings with the point: it is the sum of the two lower bounds, part~(i).
What remains to be bounded is the gradient alone, and parts~(ii) and~(iii)
say how an enclosure of it yields the two lower bounds, the feasible one from
a residual known up to a ball around an exact matrix, the frame of the rows
enclosed with it, so that no deduction for the frame arises.
Corollary~\ref{cor:certified-keylength} takes the key length from this
lemma, and Requirement~\ref{req:certified-evaluation} is how a run forms
it. Lemma~\ref{lem:certified-bracket-enclosed} says when one enclosure of
the gradient serves a reference companion known through enclosures.

\begin{lemma}\label{lem:certified-bracket}
Let $\alpha\in(1,\infty)$, let $\Decl$ be a tuple of the form of
Definition~\ref{def:declaration}, in the \tref{basic format} or the general
format, that satisfies conditions (i)--(vi) and (viii) of
Definition~\ref{def:admissible-declaration}, and let $\acSet$ satisfy
Definition~\ref{def:admissible-acceptance-set}. Let $x=(\mathbf q,J)$ with
$\mathbf q\in\acSet$ and every block of $J$ positive definite, on every
\tref{joint sector} in the general format, satisfy (H-POS) of
Theorem~\ref{thm:firstorder}, and let $(\lmoq,\lmoJ)$ be the gradient of
$O^{J}_{\alpha}$ at $x$ assembled by Corollary~\ref{cor:full-gradient},
read with Lemma~\ref{lem:fo-khyp}\,(iii) and, in the general format, on the
joint sectors as Theorem~\ref{thm:firstorder} reads it.
\begin{enumerate}[label=(\roman*),leftmargin=2.4em]
\item \emph{(The Euler form.)}
  \begin{equation}\label{eq:euler-form}
    \langle\lmoq,\mathbf q\rangle+\langle\lmoJ,J\rangle=O^{J}_{\alpha}(x).
  \end{equation}
  Consequently, for real numbers $b^{\mathrm A}$ and $b^{\mathrm F}$ with
  $b^{\mathrm A}\le\min_{\mathbf q'\in\acSet}\langle\lmoq,\mathbf q'\rangle$
  and $b^{\mathrm F}\le\min_{J'\in\Feas^{J}}\langle\lmoJ,J'\rangle$, the
  minimum over the exact $\Feas^{J}$, the bracket of
  Theorem~\ref{thm:firstorder}\,(iv) at $x$ with both pairings replaced is
  $b^{\mathrm A}+b^{\mathrm F}$, and
  \begin{equation}\label{eq:certified-bracket}
    \max\{0,\,b^{\mathrm A}+b^{\mathrm F}\}\;\le\;\hlow ,
  \end{equation}
  in exact real arithmetic.
\item \emph{(The acceptance side from an enclosure.)} If
  $\hat g\in\mathbb{R}^{\Chacc}$ and $\delta\ge0$ satisfy
  $|\lmoq(\hat c)-\hat g(\hat c)|\le\delta$ for every $\hat c\in\Chacc$, then
  $\min_{\mathbf q'\in\acSet}\langle\lmoq,\mathbf q'\rangle\ge
  \min_{\mathbf q'\in\acSet}\langle\hat g,\mathbf q'\rangle-\delta$.
\item \emph{(The feasible side from enclosed residuals.)} For each group of
  Lemma~\ref{lem:feasible-side-dual}\,(ii), an Alice sector $\mathsf t$ in
  \tref{trusted-source mode}, its rows built from the exact frame
  $W_{\mathsf t}$ of Theorem~\ref{thm:choi}, or the one group in
  \tref{untrusted-source mode}, let $Y$ and the $\Lambda_r$ be as in
  Lemma~\ref{lem:feasible-side-dual}\,(i), let $R_p$ be the residual
  \eqref{eq:dual-residual} formed with $\lmoJ$, and for every block $p$ of the
  group let $M_p$ be a Hermitian matrix, $\rho_p\ge0$ with
  $\|R_p-M_p\|_{\mathrm F}\le\rho_p$, the Frobenius norm of
  \xref{foundations/cor:least-eigenvalue-over-enclosure}, and $\ell_p$ a
  lower bound on $\lmin(M_p)$. Then the group's
  $\lambda\coloneqq\min_p(\ell_p-\rho_p)$ satisfies $\lambda\le\lmin(R_p)$
  for every $p$, and the sum over the groups of $b(\lmoJ;Y,\Lambda,\lambda)$
  of \eqref{eq:feasible-side-dual} is a number
  $b^{\mathrm F}\le\min_{J'\in\Feas^{J}}\langle\lmoJ,J'\rangle$.
\item \emph{(The \tref{dropped program}.)} In the setting of
  Corollary~\ref{cor:fo-dropped}, with a \tref{dropped set} $\mathcal{D}$,
  (i)--(iii) hold with $O^{J,\mathcal{D}}_{\alpha}$ in place of
  $O^{J}_{\alpha}$, (H-POS$^{\mathcal{D}}$) in place of (H-POS), its gradient
  of (H-NOPERT$^{\mathcal{D}}$) in place of $(\lmoq,\lmoJ)$, the kept pairs in
  place of the pairs that are not key-inert, and
  $\underline{h}^{\mathcal{D}}_{\alpha}$ in place of $\hlow$; and
  $\underline{h}^{\mathcal{D}}_{\alpha}\le\hlow$.
\end{enumerate}
\end{lemma}

\begin{proof}
\emph{Idea.} Homogeneity makes the derivative along the point equal to the
value, and the gradient's formulas give the same derivative as the pairing
with the point; the two lower bounds then carry the whole bracket. An
entrywise enclosure of the acceptance-side gradient moves its minimum over a
set of distributions by at most the radius, and a ball around the residual
moves its least eigenvalue by at most the ball's radius.

\emph{(i), Step 1 (the point is admissible to the theorem).} Every block of
$J$ is positive semidefinite, and each residual of
Definition~\ref{def:relaxed-feasible-set} at $J$, of the coupling condition,
of the normalisation and of every declared row, is a finite number, so $J$
lies in $\Feas^{J}(\varepsilon)$ for every $\varepsilon$ at least each of
them, the normalisation's divided by the dimension factor its clause
carries; with
$\mathbf q\in\acSet$, $x\in\acSet\times\Feas^{J}(\varepsilon)$. Every block
being positive definite, every key-active branch operator attains its
\tref{maximal branch support} (Lemma~\ref{lem:branch-support}\,(i), per joint
sector in the general format), which is (H-SUPP); and $(\lmoq,\lmoJ)$ is the
gradient at $x$ itself, which is (H-NOPERT).

\emph{Step 2 (the value along the ray).} By (H-POS) and
Lemma~\ref{lem:fo-khyp}\,(i), $q(\hat c)>0$ exactly on $\Chacc$, and
$P^{J}_{\hat c}(J)>0$ there, $K^{J}_{\hat c}(J)>0$ on $\Chkey\cap\Chacc$, so
every term of $O^{J}_{\alpha}(x)$ indexed in $\Chacc$ is real, and every
term indexed outside it has $q(\hat c)=0$ and is $0$ by
\xref{foundations/def:kl}, the second arguments being non-negative by the
cone clause of Theorem~\ref{thm:branch-formula}. So $O^{J}_{\alpha}(x)$ is
real. By Lemma~\ref{lem:objective-homogeneous}\,(iii),
$O^{J}_{\alpha}((1+\tau)x)=(1+\tau)\,O^{J}_{\alpha}(x)$ for every
$\tau>-1$, so $\tau\mapsto O^{J}_{\alpha}(x+\tau x)$ is differentiable at
$\tau=0$ with derivative $O^{J}_{\alpha}(x)$.

\emph{Step 3 (the same derivative from the gradient).} The terms indexed
outside $\Chacc$ vanish at every $(1+\tau)x$, $\tau>-1$, by Step~2 applied to
the scaled point, $(1+\tau)q(\hat c)=0$ there; they contribute $0$ to the
derivative, and the pairing $\langle\lmoq,\mathbf q\rangle$ of
Definition~\ref{def:first-order-setting} and the sums of \eqref{eq:grad-J}
read with Lemma~\ref{lem:fo-khyp}\,(iii) omit them. For $\hat c\in\Chacc$
the term is a smooth function of $(q(\hat c),P^{J}_{\hat c},K^{J}_{\hat c})$
near the point, these being positive at $x$ and scaling with $1+\tau$
(Lemma~\ref{lem:objective-homogeneous}\,(i)). Along $\tau$,
$q(\hat c)$ and $P^{J}_{\hat c}$ move linearly, with derivatives $q(\hat c)$
and $P^{J}_{\hat c}(J)=\sum_{s'}\tr[\Gamma^{(s')}_{\hat c}J_{s'}]$
(\eqref{eq:grad-pieces}). Each deficit of a key-active pair is
$f^{J}_{\hat c,s}(J_s+\tau J_s)=T(X+\tau X)-\tr(X+\tau X)$ with
$X\coloneqq\cG^{J}_{\hat c,s}(J_s)$, the \tref{branch map} being linear; the
direction $X$ is Hermitian with support $\operatorname{supp}X$, so
Theorem~\ref{thm:gradient}\,(i) gives the derivative
$\tr[\nabla T(X)\,X]-\tr X=\tr\big[(\nabla T(X)-\id)\,\cG^{J}_{\hat c,s}(J_s)\big]
=\tr\big[\chi^{\dagger}_{\Qsys_s}\big(\cG^{\dagger}_{\hat c,s}[\nabla T(X)-\id]\big)
J_s\big]$, which is $\tr[\Delta^{(s)}_{\hat c}J_s]$ of \eqref{eq:grad-pieces};
the deficits of key-inert pairs are identically $0$
(Definition~\ref{def:keyactive}, Lemma~\ref{lem:joint-sectors}\,(iii)).
The chain rule through $u\log_2\frac uv$ at $u,v>0$, exactly as in the proof
of Corollary~\ref{cor:full-gradient}, then gives the derivative of the term as
its coefficients of \eqref{eq:grad-q} and \eqref{eq:grad-J} paired with
$q(\hat c)$ and with the blocks of $J$, and summing over $\Chacc$, the
derivative of $\tau\mapsto O^{J}_{\alpha}(x+\tau x)$ at $0$ is
$\langle\lmoq,\mathbf q\rangle+\langle\lmoJ,J\rangle$. With Step~2 this is
\eqref{eq:euler-form}.

\emph{Step 4 (the bracket).} By Step~1, Theorem~\ref{thm:firstorder}\,(iiib)
with its part~(iv) applies to the one point $x$ with both pairings replaced,
where $\Decl$ satisfies condition~(vii) as well; its bracket is
$O^{J}_{\alpha}(x)+b^{\mathrm A}-\langle\lmoq,\mathbf q\rangle
+b^{\mathrm F}-\langle\lmoJ,J\rangle=b^{\mathrm A}+b^{\mathrm F}$ by
\eqref{eq:euler-form}, and the theorem gives \eqref{eq:certified-bracket}.
Where condition~(vii) fails, Theorem~\ref{thm:firstorder}\,(v) gives the same,
$\hlow=+\infty$.

\emph{(ii).} Let $\mathbf q'\in\acSet$, a non-empty set as it contains
$\mathbf q$. It is a probability vector vanishing outside $\Chacc$
(Lemma~\ref{lem:fo-khyp}\,(i)), so $q'(\hat c)\ge0$ and
$\sum_{\hat c\in\Chacc}q'(\hat c)=1$. Hence
$\langle\lmoq,\mathbf q'\rangle=\langle\hat g,\mathbf q'\rangle
+\sum_{\hat c\in\Chacc}(\lmoq(\hat c)-\hat g(\hat c))q'(\hat c)
\ge\langle\hat g,\mathbf q'\rangle-\delta
\ge\min_{\mathbf q''\in\acSet}\langle\hat g,\mathbf q''\rangle-\delta$, the
minimum attained, $\acSet$ being compact
(Definition~\ref{def:admissible-acceptance-set}) and the map linear. The
minimum over $\mathbf q'$ of the left-hand side is therefore at least the
right-hand side.

\emph{(iii).} $R_p$ is Hermitian, $\lmoJ$, $Y\otimes\id$ and every
$L^{*}_{r,p}(\Lambda_r)$ being so
(Lemma~\ref{lem:feasible-side-dual}\,(i), (ii)). By
\xref{foundations/cor:least-eigenvalue-over-enclosure}\,(i),
$\lmin(R_p)\ge\lmin(M_p)-\rho_p\ge\ell_p-\rho_p\ge\lambda$. So $\lambda$ is
admissible in Lemma~\ref{lem:feasible-side-dual}\,(i) for the group, whose
rows, built from the exact frame, cut out the factor of the exact
$\Feas^{J}$ (part~(ii)(a) of that lemma; part~(ii)(b) in untrusted-source
mode), and the sum over the groups bounds $\langle\lmoJ,J'\rangle$ from
below for every $J'\in\Feas^{J}$, which is
\eqref{eq:feasible-side-sum}.

\emph{(iv).} $O^{J,\mathcal{D}}_{\alpha}$ is positively homogeneous of degree
one (Lemma~\ref{lem:objective-homogeneous}\,(iii)), and Steps~2 and~3 run
unchanged with $K^{J,\mathcal{D}}_{\hat c}$ in place of $K^{J}_{\hat c}$
and the deficits of the kept pairs only, the dropped ones being absent from
\eqref{eq:KD}; (H-SUPP$^{\mathcal{D}}$) holds at a positive definite block as
in Step~1. Step~4 runs with Corollary~\ref{cor:fo-dropped}, whose
conclusion \eqref{eq:hlb-dropped} admits the lower bounds of
Theorem~\ref{thm:firstorder}\,(iv), in place of the theorem. Parts (ii) and
(iii) do not read the objective. Finally
$\underline{h}^{\mathcal{D}}_{\alpha}\le\hlow$ is
Proposition~\ref{prop:drop-bound}.
\end{proof}
